\subsection{B7. Abgleich des zuletzt eingesandten Berichts
v1.6.3}\label{b7.-abgleich-des-zuletzt-eingesandten-berichts-v1.6.3}

Der sechste Text dieser Runde, „Operationssatz, nativer Grundzustand und
Feldwörterbuch``, enthält überwiegend bereits berücksichtigte Befunde:
die Kommutantenleiter, das Lochbild, die Normen der ersten Stufen, den
nichtverschwindenden Seitenzweig und den verschwindenden
Weyl-Skalarkanal. Seine zweite vorgeschlagene Folgeaufgabe wurde
inzwischen wesentlich weitergeführt: Alle vier Zwei-Boson-Singuletts
sind konstruiert und der richtige H-Lanczos-Anfang ist fünf Glieder weit
bestimmt.

Vier Präzisierungen verhindern eine Übernahme veralteter Schlüsse:

\begin{enumerate}
\def\labelenumi{\arabic{enumi}.}
\tightlist
\item
  \textbf{Sieben und sechzehn sind verschiedene Kommutanten.} Sieben
  gilt für volle innere Gruppe zusammen mit X und Nb; sechzehn für die
  Gruppe allein. Die drei hellen inneren Typen kommen auf beiden Seiten
  der Umwandlung vor. Einzelne Tabellenzeilen sind multiplizitätsfrei,
  der ganze N=3-Sektor aber nicht. Der Schluss vom nichtskalaren Clock
  auf eine äußere, nicht zusammenhängende Symmetrie ist weiterhin
  falsch.
\item
  \textbf{Eine kleine Zweignorm beweist keine globale Konvergenz.} Der
  genannte Defekt 2,90 mal 10 hoch -5 betrifft eine bestimmte
  Kontraktion. Die höhere Bosonstufe ist weiterhin gekoppelt. „Numerisch
  fast allein durch die Normen lösbar`` und „vernachlässigbar`` sind
  ohne Restschranke zu stark. Auch Zh plus Zadd gleich eins ist eine
  Summenregel, keine Bestimmung sämtlicher Spektrallinien.
  Ritz-Erwartungswerte bleiben Erwartungswerte des Ritz-Zustands, nicht
  automatisch von Omega.
\item
  \textbf{Der Feldschluss ist ansatzabhängig.} Der angegebene Fünfzyklus
  (16,45,0,30,37) wurde frisch am W überprüft. Damit ist eine einzige
  binäre Chiralitätsbelegung der primitiven Labels mit entgegengesetzter
  Händigkeit auf jeder Kante unmöglich. Das verbietet kein allgemeines
  Dirac-Wörterbuch mit erweiterten Trägern, Ableitungen oder
  zusammengesetzten Feldern. Der Tensorfeldkanal ist unter dem
  angegebenen engen Ansatz sinnvoll, nicht schon der universell einzige
  physikalische Feldtyp.
\item
  \textbf{Der gespeicherte Rechnerstand ist älter als der Text.} Die
  lokal zugehörige Datei \codeword{native\_ground\_state.json} hat einen
  anderen Checkerhash als das aktuelle Programm. Ihr gespeichertes
  w2-Normquadrat ist um 229 Quadrat gegenüber der physikalischen Norm
  skaliert. Das aktuelle Programm enthält bereits die richtige
  Entskalierung; die alte Datei ist also kein aktueller Replaynachweis.
  Sie meldet außerdem \codeword{norms\_by\_traces\_used=false} und
  Ritz-Tiefe drei. Die im Text erwähnte Datei
  \codeword{norms\_by\_traces.json} und das dortige Replaymanifest waren
  bei dieser Prüfung nicht vorhanden. Diese Beobachtung schließt ein
  Resultat an anderem Ort nicht aus, liefert hier aber keines.
\end{enumerate}

Daher wird insbesondere die Behauptung, nu4 sei in dieser Quelle auf
zwei unabhängigen Wegen vollständig reproduziert, nicht übernommen. Der
Text selbst bezeichnet nu4 zugleich als „nur Spurnetzwerk`` und
beschreibt den Brute-Force-Abgleich nur bis nu3. Der frühere gepinnte
exakte nu4-Eingang unserer eigenen Rechnung bleibt davon getrennt und
wurde nicht als neue Enumeration ausgegeben. Ein fertig berechnetes nu5
wird aus dem vorhandenen Programm mit \codeword{MAX\_N=5} nicht
abgeleitet; eine Einstellung ist kein Ergebnis.

\textbf{Die interessante verbleibende Richtung ist die
Spurnetzwerk-Methode.} Sie kann Normen durch Kontraktionen des
vorhandenen Tensors berechnen, statt alle Vielteilchenamplituden zu
speichern. Zur fundamentalen Vereinfachung wird sie erst dann, wenn sie
auch gemischte H-Momente, Seitenzweige und kontrollierte
Mehrzeitantworten liefert. Der konkrete nächste Test ist daher kein
bloßer weiterer Normwert: Eine zweite kleine Kontraktionsrechnung soll
die schon exakt bekannten zehn H-Momente reproduzieren und dann
mindestens ein noch nicht bekanntes gemischtes Matrixelement samt
Fehlerkontrolle bestimmen. Diese Methode lässt sich mit der
Schattenkarte verbinden, ersetzt aber keine native Operation.

Das Prüfpaket bewahrt neben dem Text die beiden widersprechenden
Programmstempel und das vorhandene Spurnetzwerkprogramm. Diese fremden
Dateien wurden nicht verändert oder als vollständig neu ausgeführt
ausgegeben.

\begin{center}\rule{0.5\linewidth}{0.5pt}\end{center}

\section{Historischer Anhang: vollständige Fassung
v1.6.6}\label{historischer-anhang-vollstuxe4ndige-fassung-v1.6.6}

Die folgende vollständige Fassung wird unverändert bewahrt. Maßgeblich
für den aktuellen Status sind die vorangestellten Kapitel v1.6.7. Ältere
Vorhaben, Modellannahmen und Zahlen erhalten dadurch keinen neuen
Beweisstatus; die kleinen Energieverschärfungen und die neuen Singulett-
und Bilinearergebnisse stehen oben.

\section{TFPT / Universalraum: Clock, gemeinsame Quelle und
tatsächlicher
Transport}\label{tfpt-universalraum-clock-gemeinsame-quelle-und-tatsuxe4chlicher-transport}

\subsection{Konsolidierung und eigene Forschungsfortsetzung v1.6.6 - 15.
September
2026}\label{konsolidierung-und-eigene-forschungsfortsetzung-v1.6.6---15.-september-2026}

\textbf{Ergebnis:} Der dokumentierte endliche Clock ist jetzt
ausdrücklich als Element der vorhandenen inneren Spin(10)-Symmetrie
konstruiert. Dadurch lässt sich die bislang offene gemeinsame
Clock-/Casimir-Auslesung exakt bestimmen. Daneben wurde ein kleiner,
echter Vierzustandskanal der ursprünglichen Paarwechselwirkung gefunden,
der nach Zulassung eines Bosonmischers eine Fermionenmarke überträgt.
Die Herkunft dieses Mischers ist nicht bewiesen.

Der neue Universalraum-Text enthält eine sinnvolle Suchrichtung -
gemeinsame Quelle statt vorab unabhängiger Banken -, aber Überlappung
allein erzeugt weder Transport noch Raumzeit, Eichkrümmung oder die
arithmetische Spur. Mehrere seiner vermeintlich automatischen Übergänge
werden hier präzisiert.

Diese Revision integriert \textbf{drei neue Nutzertexte}, die
unabhängigen parallelen Prüfungen und eigene neue Herleitungen. Die
vollständige frühere Konsolidierung v1.6.5 einschließlich v1.6.4 bleibt
im historischen Anhang erhalten. Das ist eine aktualisierte
Forschungsfassung in Markdown mit kurzem Änderungsdokument und einfacher
Erklärung; keine behauptete neue PDF-/Web-Publikation oder vollständige
Theory of Everything.

\subsubsection{Leseschlüssel}\label{leseschluxfcssel}

\begin{itemize}
\tightlist
\item
  \textbf{Exakt:} Identität oder mathematische Folgerung im angegebenen
  Modell.
\item
  \textbf{Numerisch:} berechneter Näherungswert, ohne zertifizierte
  Einschließung.
\item
  \textbf{Bedingt:} exakter Satz erst nach ausdrücklich zusätzlicher
  Ressource.
\item
  \textbf{Offen:} fehlende Herleitung oder physikalische Identifikation.
\end{itemize}

Eine grüne Prüfung ersetzt weder die Voraussetzungen eines Satzes noch
eine ausführbare Operation. Assoziative Algebradimension, dynamische
Kontrollierbarkeit und physische Verfügbarkeit werden getrennt geführt.

\subsection{1. Die unveränderte gemeinsame
Grundlage}\label{die-unveruxe4nderte-gemeinsame-grundlage}

Mit 64 Fermionmoden, 60 Bosonmoden und demselben gepinnten Tensor gilt

\[\adjustbox{max width=\linewidth}{$\displaystyle 
H_{\rm nat}=\Delta N_b+g\sum_A(b_A^\dagger P_A+P_A^\dagger b_A),
\qquad P_A=\sum_{i<j}W_{A,ij}f_jf_i,
\qquad N=N_f+2N_b.
$}\]

W besitzt 480 Einträge ±1, acht disjunkte Paare je Zeile, und
\codeword{WW†=8 I60}. Die Fermionmarkierung ist \codeword{16⊗4}, die
Bosonmarkierung \codeword{10⊗6}. Spin(10)×SU(4) ist zunächst eine
\textbf{innere} Quellsymmetrie, keine bereits hergeleitete
Raumzeit-Lorentzgruppe.

Der frühere native Satz bleibt unverändert maßgeblich: Für
Δ\textgreater0 und \codeword{0<|g|/Δ≤1/20} hat dieser μ=0-Vertrag einen
eindeutigen globalen Grundzustand Ω bei N=64. Am Prüfpunkt
\codeword{g/Δ=1/20} gelten unter anderem

\[\adjustbox{max width=\linewidth}{$\displaystyle 
-1.158089\Delta<E_0<-1.129636\Delta,
\qquad 0.842846<\langle N_b\rangle_\Omega<1.245656.
$}\]

Die isolierte niedrige Entnahmelinie im N=63-Sektor hat weiterhin die
bewiesenen Schranken aus v1.6.4/v1.6.5:

\[\adjustbox{max width=\linewidth}{$\displaystyle 
0.007737\Delta<\epsilon<0.039079764\Delta,
\qquad
Z_{\rm low}>\frac{40912436089}{46487375000}>0.88007628.
$}\]

Z bezieht sich auf das gesamte normierte Fermionspektralgewicht pro
Mode. Die gesamte Entnahmenorm ist dagegen
\codeword{ν=<f\_r†f\_r>=1-<Nb>/32}. Die beiden Größen sind nicht
austauschbar. Der Pol ist eine wohldefinierte isolierte Eigenlinie
dieses endlichen Modenmodells; seine relativistische Teilchen-, Massen-
oder räumliche Propagationsdeutung ist eine weitere Aufgabe.

Der vollständige bedingte Zwei-Banken-Satz aus v1.6.5 bleibt erhalten:
mit deklariertem Fermionlink und Rotorvertrag gelingt der reine niedrige
Transfer mit Wahrscheinlichkeit \textgreater99,267 \%, die unfiltrierte
ursprüngliche Entnahme mit \textgreater89,726 \%. Der Link war und ist
\textbf{zusätzlich vorausgesetzt}. Die neuen Vierzustandszahlen unten
gehören zu einem anderen Vertrag und ersetzen diese Ergebnisse nicht.

\subsection{2. Quellenprüfung der beiden neuen
Ergebnisrunden}\label{quellenpruxfcfung-der-beiden-neuen-ergebnisrunden}

\subsubsection{2.1 Bestätigte Zerlegung, korrigierte
Operationsaussagen}\label{bestuxe4tigte-zerlegung-korrigierte-operationsaussagen}

Die 944, 31 und 152 mitgelieferten Prüfbedingungen der Symmetrie-,
Verfügbarkeits- und Lorentzprogramme sind normal und optimiert
byteidentisch reproduziert. Darunter befinden sich neun wörtliche
\codeword{need(True, ...)}-Guards mit theoretischen Folgerungen. Sie
zählen als ausgeführte Guards, nicht als neun unabhängige Beweise. Die
Clock-Normalisatorbehauptung gehörte dazu; Abschnitt 3 ersetzt sie durch
eine explizite Konstruktion.

Im gesamten N=3-Raum gilt die Zerlegung:

{\def\LTcaptype{none} % do not increment counter
\begin{longtable}[]{@{}
  >{\raggedright\arraybackslash}p{(\linewidth - 8\tabcolsep) * \real{0.1579}}
  >{\raggedleft\arraybackslash}p{(\linewidth - 8\tabcolsep) * \real{0.2105}}
  >{\raggedleft\arraybackslash}p{(\linewidth - 8\tabcolsep) * \real{0.2105}}
  >{\raggedleft\arraybackslash}p{(\linewidth - 8\tabcolsep) * \real{0.2105}}
  >{\raggedleft\arraybackslash}p{(\linewidth - 8\tabcolsep) * \real{0.2105}}@{}}
\toprule\noalign{}
\begin{minipage}[b]{\linewidth}\raggedright
Typ
\end{minipage} & \begin{minipage}[b]{\linewidth}\raggedleft
irreduzible Dimension
\end{minipage} & \begin{minipage}[b]{\linewidth}\raggedleft
Multiplizität
\end{minipage} & \begin{minipage}[b]{\linewidth}\raggedleft
Spin-Casimir
\end{minipage} & \begin{minipage}[b]{\linewidth}\raggedleft
Farb-Casimir
\end{minipage} \\
\midrule\noalign{}
\endhead
\bottomrule\noalign{}
\endlastfoot
(1200,20) & 24000 & 1 & 141/4 & 39/4 \\
(560,20′) & 11200 & 1 & \textbf{117/4} & 63/4 \\
(672,4̄) & 2688 & 1 & 165/4 & 15/4 \\
(144,20) & 2880 & 2 & 85/4 & 39/4 \\
(144,4̄) & 576 & 2 & 85/4 & 15/4 \\
(16̄,20) & 320 & 2 & 45/4 & 39/4 \\
(16̄,4̄) & 64 & 1 & 45/4 & 15/4 \\
\end{longtable}
}

Die gewichtete Gesamtdimension ist 45504. Nur die getrennten Teile Λ³64
und 60⊗64 sind jeweils multiplizitätsfrei, nicht ihre direkte Summe. Die
drei dunklen Typen besitzen alle Gesamtcasimir 45, werden aber bereits
durch \textbf{einen} getrennten Casimir unterschieden.

Eine unabhängige rationale Rechnung auf den kleinen Multiplizitätsräumen
liefert für die jeweils \textbf{gewährten} Operationen:

{\def\LTcaptype{none} % do not increment counter
\begin{longtable}[]{@{}lr@{}}
\toprule\noalign{}
Operationsvertrag & komplexe assoziative Algebradimension \\
\midrule\noalign{}
\endhead
\bottomrule\noalign{}
\endlastfoot
Ein fixes H=Nb+X/20 & 8 \\
X und Nb unabhängig & 14 \\
Zusätzlich ein dunkler Isotypieprojektor & 15 \\
Zusätzlich ein getrennter Casimir & 16 \\
Zusätzlich nur Gesamtcasimir & 14 \\
\end{longtable}
}

Somit gibt es keine universelle binäre Frage „alle Symmetriegriffe oder
gar kein Fortschritt``. Die Dimension 15 ist ein konkretes
Gegenbeispiel. Die vollständige Gruppen-Generatoralgebra hat mit X,Nb
die Dimension 743583744 und Kommutantdimension 7. Sie ist unter
Gruppenkonjugation \textbf{stabil}, aber nicht punktweise invariant.
Sieben ist kein universeller Boden für beliebige symmetriestabile
Algebren: End(H) wäre ebenfalls stabil und hätte einen eindimensionalen
Kommutanten. Keine dieser Zahlen beweist, dass die zugehörigen
Operationen im Compiler ausführbar sind.

\subsubsection{2.2 Der Grundzustandsbericht enthält einen
Versionskonflikt}\label{der-grundzustandsbericht-enthuxe4lt-einen-versionskonflikt}

Der neue Text hat die Nichtschlussnorm bereits korrigiert, das
danebenliegende \codeword{native\_ground\_state.json} stammt jedoch von
einer älteren Codefassung. Sein als PASS markierter Wert ist um
\codeword{229²} zu groß; die dortige relative orthogonale Norm ist sogar
größer als eins. Der daraus berechnete Ritzwert −1,18826198 wird
\textbf{nicht} übernommen.

Eine neue direkte ganzzahlige Kontraktion umgeht die 15,25 Millionen
Komponenten von v₃. Sie berechnet \codeword{u=T T†v₂} auf 293280
nichtverschwindenden Komponenten und bestätigt exakt

\[\adjustbox{max width=\linewidth}{$\displaystyle 
\|u\|^2=752194252800,\quad
\langle v_2,u\rangle=575078400,\quad
w_2=u-\frac{299520}{229}v_2,
\quad\|w_2\|^2=\frac{5001523200}{229}>0.
$}\]

Der eine Vektor je Bosonlage schließt daher nicht. w₂ ist ein zweiter
Singulettvektor derselben Lage. Für die kleine vier- bzw.
fünfdimensionale Ritzmatrix ergeben sich bei g/Δ=1/20 \textbf{numerisch}

\[\adjustbox{max width=\linewidth}{$\displaystyle 
E_{R,4}/\Delta=-1.0942308394409612,
\qquad E_{R,5}/\Delta=-1.0942318710142458.
$}\]

Die zusätzliche w₂-Richtung ändert diese Näherung tatsächlich erst um
etwa 10⁻⁶. Das bedeutet \textbf{keine} Konvergenz zum wahren
Grundzustand: Beide Werte liegen noch über dessen bereits bewiesener
Obergrenze.

Die neue einfache Restformel zeigt die Ursache. Für den normierten
Vier-Vektor-Ritzeigenzustand mit letzter Komponente c₃ gilt

\[\adjustbox{max width=\linewidth}{$\displaystyle 
\|(H-E_R)\psi_3\|^2
=g^2|c_3|^2\frac{\nu_4+\|w_2\|^2}{\nu_3}.
$}\]

Mit dem früher geprüften vierten Moment ist die Restnorm numerisch rund
0,373063 Δ. Der kleine w₂-Zweig trägt weniger als 23 Millionstel des
quadrierten Restes; die ausgelassene v₄-Richtung dominiert. Die neuen
Zahlen \codeword{Zh≈0,9737} und \codeword{εh≈0,031Δ} gehören ebenfalls
zur Ritz-Näherung: Gesamtentnahmenorm und erstes Moment, nicht neu
bestimmte Polgrößen auf Ω.

Die Suche nach \codeword{E<−1,2Δ} oder \codeword{E<−1,1625Δ} wäre sogar
mit der bestehenden unteren Schranke unvereinbar. Die schwächere neue
Beweismethode öffnet den bereits geschlossenen nativen Grundzustandssatz
nicht wieder.

\subsection{3. Neuer exakter Anschluss: Der Clock liegt bereits in
Spin(10)}\label{neuer-exakter-anschluss-der-clock-liegt-bereits-in-spin10}

\subsubsection{3.1 Ausgangspunkt ist der tatsächliche
Quell-Clock}\label{ausgangspunkt-ist-der-tatsuxe4chliche-quell-clock}

Der gepinnte endliche Compilerpräfix liefert auf fünf Hilfsachsen

\[\adjustbox{max width=\linewidth}{$\displaystyle 
p=(0\mapsto2,\ 1\mapsto0,\ 2\mapsto1,\ 3\mapsto4,\ 4\mapsto3),
\qquad\det p=-1.
$}\]

Seine vorzeichenrichtige Exteriorhebung auf dem geraden
5-Moden-Hilfsfockraum ist Γ(p)\textbar even. Auf den nativen Fermionen
ist \codeword{GF=Γ(p)|even⊗I4}; auf den Bosonen ist
\codeword{GB=(-P10)⊗I6}, wobei P10 beide Fünfergruppen permutiert. Diese
Darstellung wird aus dem vorhandenen Clockpräfix erneut extrahiert,
nicht aufgrund passender Eigenwerte ausgewählt.

